\section{单圈与两圈匹配核}\label{sec:lamet-oneloop}

显式的定阶结果在所述的阶上是严格的；它们检验全阶匹配定理（\cref{sec:lattice-matching}），也是格点分析的输入。单圈图见\cref{fig:lamet-oneloop}，记号见\cref{tab:lamet-conventions}（$D=[0,1]$ 上的加号分布简记为 $[g]_+$）。

%%% FIGBLOCK fig:lamet-oneloop

\begin{proposition}[准 PDF 匹配的单圈结果 \tagR]\label{prop:lamet-oneloop}
（$R_{\mathrm{PT}}$：在重整化微扰论中对红外正规化的部分子态严格。）在单圈、味非单态、算符按 $\MSbar$ 重整化（单圈抵消项在此计算，它就是\cref{prop:lattice-renormalizability}(d) 的单圈内容）、部分子态以维数正规化为红外正规子时：
\begin{enumerate}[label=(\alph*),leftmargin=2em,itemsep=2pt,topsep=2pt]
\item 坐标空间核（$\Gamma=\gamma^0$：$\sigma=0$；$\Gamma=\gamma^z$：$\sigma=1$）
\begin{equation}\label{eq:lamet-coord-kernel}
\begin{split}
\mathcal C(\alpha,L_z)={}&\Big[1+a\Big(\frac32L_z+\frac32\Big)\Big]\delta(1-\alpha)\\
&+a\Big\{\Big[\frac{1+\alpha^2}{1-\alpha}\Big]_+(-L_z-1)-\Big[\frac{4\ln(1-\alpha)}{1-\alpha}\Big]_++2(1+\sigma)(1-\alpha)\Big\}\theta(\alpha)\theta(1-\alpha),
\end{split}
\end{equation}% VF-LAM-4
其矩为 $C_n=\int\dd\alpha\,\alpha^n\mathcal C=1+a\big[\big(\frac{3+2n}{2+3n+n^2}+2H_n\big)L_z+\frac{5+2n}{2+3n+n^2}+2(1-H_n)H_n-2H^{(2)}_n\big]$（$\gamma^z$ 另加 $\frac{2a}{2+3n+n^2}$；$H_n=\sum_{i\le n}\frac1i$，$H^{(2)}_n=\sum_{i\le n}\frac1{i^2}$），特别地 $C_0=1+a(\frac32L_z+\frac52)$（$\gamma^z$：$\frac72$）；重整化群方程 $\frac{\dd\mathcal C}{\dd\ln\mu^2}=a[\frac32\delta(1-\alpha)-B(\alpha)]$，$B=[\frac{1+\alpha^2}{1-\alpha}]_+$。
\item 比值核 $\mathcal C^{\rm ratio}\equiv\mathcal C/C_0$：在单圈 $\mathcal C^{\rm ratio}=\mathcal C-\delta(1-\alpha)(C_0-1)=\delta(1-\alpha)+a\big\{B(-L_z-1)-[\frac{4\ln(1-\alpha)}{1-\alpha}]_++(1+\sigma)[2(1-\alpha)]_+\big\}$，它的单圈部分全是加号分布；它与 Radyushkin 的约化 ITD 关系相同（他的 $\nu=-\lambda$）。% VF-LAM-5 (the O(a) part of VF-LAM-5 is written with the tree-level delta added)
\item 非物理区的准 PDF 即\cref{ex:lattice-quasi-tree}(b) 的\cref{eq:lamet-unphysical}。
\item 动量空间核（在固定部分子动量下分别是 (a)(b) 的分布意义下的 Fourier 变换；$\xi=\pm\infty$ 处的 $\delta$ 函数已略去，它们对 $f(y)\sim y^{-1+a'}$、$a'>0$ 不起作用）：
\begin{equation}\label{eq:lamet-CMS}
\begin{split}
C^{\MSbar}(\xi)={}&\delta(1-\xi)\Big[1+a\Big(\frac32\ell_p+\frac52\Big)\Big]\\
&+a\begin{cases}\big[\frac{1+\xi^2}{1-\xi}\ln\frac\xi{\xi-1}+1+\frac3{2\xi}\big]^{[1,\infty]}_{+(1)}-\frac3{2\xi},&\xi>1\\[3pt]\big[\frac{1+\xi^2}{1-\xi}\big(-\ell_p+\ln(\xi(1-\xi))\big)-\frac{\xi(1+\xi)}{1-\xi}\big]^{[0,1]}_{+(1)},&0<\xi<1\\[3pt]\big[-\frac{1+\xi^2}{1-\xi}\ln\frac{-\xi}{1-\xi}-1+\frac3{2(1-\xi)}\big]^{[-\infty,0]}_{+(1)}-\frac3{2(1-\xi)},&\xi<0\end{cases},
\end{split}
\end{equation}
\begin{equation}\label{eq:lamet-Cratio}
C^{\rm ratio}(\xi)=\delta(1-\xi)+a\begin{cases}\big[\frac{1+\xi^2}{1-\xi}\ln\frac\xi{\xi-1}+1-\frac3{2(1-\xi)}\big]^{[1,\infty]}_{+(1)},&\xi>1\\[3pt]\big[\frac{1+\xi^2}{1-\xi}\big(-\ell_p+\ln(\xi(1-\xi))-1\big)+1+\frac3{2(1-\xi)}\big]^{[0,1]}_{+(1)},&0<\xi<1\\[3pt]\big[-\frac{1+\xi^2}{1-\xi}\ln\frac{-\xi}{1-\xi}-1+\frac3{2(1-\xi)}\big]^{[-\infty,0]}_{+(1)},&\xi<0\end{cases}
\end{equation}% VF-LAM-8
对 $\gamma^z$，$\MSbar$ 核加 $2a(1-\xi)\theta(0<\xi<1)$，比值核加 $a[2(1-\xi)]^{[0,1]}_{+(1)}$。
\item 端点结构（第~23~章阈值因子化的输入）：$0<\xi<1$ 时 $C^{\MSbar}\supset\frac{2a}{1-\xi}\big(L_\zeta+\ln(1-\xi)-1\big)$，$\xi>1$ 时 $C^{\MSbar}\supset\frac{2a\ln(\xi-1)}{\xi-1}$，合起来是
$a(L_\zeta-1)\big[\mathcal P\frac1{1-\xi}+\mathcal P\frac1{|1-\xi|}\big]+a\,\mathcal P\frac{\ln(1-\xi)^2}{|1-\xi|}$，其中 $L_\zeta=\ln(4p_z^2/\mu^2)$，$\mathcal P$ 表示 $\xi=1$ 处的减除处方。
\item 动量重整化群：$P^z\frac{\dd\tilde q}{\dd P^z}=\frac{\as C_F}\pi\big[\big[\frac{1+y^2}{1-y}\big]_+-\frac32\delta(1-y)\big]\otimes f+\order{\as^2}$。
\end{enumerate}
\uses{def:lattice-quasi, def:lamet-schemes, def:prelim-msbar, prop:pdf-partonic-dr, prop:pdf-lo-kernels, def:prelim-plus, ex:lattice-quasi-tree}
\end{proposition}

\begin{proof}
(c)(e)(f) 的证明是完整的；(a)(b)(d) 只给出计算的结构与关键恒等式（概要），完整的单圈积分见\cite[第~III~节与附录~B]{IzubuchiEtAl2018}与\cite{Radyushkin:2018cvn}。

(a) 夸克靶的裸单圈 ITD（各图见\cref{ex:lattice-quasi-tree}与\cite[第~III~节]{IzubuchiEtAl2018}）等于 $a(\mu|z|/2)^{2\varepsilon}\ee^{\varepsilon\gamma_E}$ 乘以：紫外极点 $\frac32\frac1{\varepsilon_{\rm UV}}\ee^{\ii\lambda}$（\cref{prop:lattice-renormalizability}(d) 的分解），加上只含红外极点与有限部分的项。按 $\MSbar$ 减去紫外极点后，红外极点正好是 $\MSbar$ 部分子分布 $f_{q/q}=\delta(1-y)-\frac{\as}{2\pi}\frac1{\varepsilon}P_{qq}$（\cref{prop:pdf-partonic-dr}，\tagR，$R_{\mathrm{PT}}$；$P_{qq}=C_F B$，\cref{prop:pdf-lo-kernels}）的 Ioffe 时间变换，二者相减即得\cref{eq:lamet-coord-kernel}；这是\cite{IzubuchiEtAl2018,Radyushkin:2018cvn,Zhang:2018ggy}的计算。矩由逐项积分得出：$\int_0^1\alpha^nB=-\frac{3+2n}{2+3n+n^2}-2H_n+\frac32$、$\int_0^1[\frac{4\ln(1-\alpha)}{1-\alpha}]_+\alpha^n=2(H_n^2+H_n^{(2)})$ 等（我们对 $n\le12$ 逐一核对了闭式）。因 $\dd L_z/\dd\ln\mu^2=1$，\cref{eq:lamet-coord-kernel}给出重整化群方程；它与 DGLAP $\frac{\dd I}{\dd\ln\mu^2}=a\,B\otimes I$ 合起来给出 $\frac{\dd\tilde h}{\dd\ln\mu^2}=\frac32a\,\tilde h+\order{a^2}$，正是\cref{prop:lattice-renormalizability}(d) 的算符反常维度。

(b) 约化 ITD $\tilde h(\lambda,z^2)/\tilde h(0,z^2)$ 的核是 $\mathcal C/C_0$（部分子靶 $\tilde h(0,z^2)=C_0$），单圈即 $\mathcal C-\delta(1-\alpha)(C_0-1)$。(a) 的 $\delta$ 项系数 $a(\frac32L_z+\frac32)$ 减去 $C_0-1=a(\frac32L_z+\frac52+\sigma)$ 得 $-a(1+\sigma)$，与 $2(1+\sigma)(1-\alpha)$ 合成 $(1+\sigma)[2(1-\alpha)]_+$（因 $\int_0^12(1-\alpha)=1$）。Radyushkin 的关系\cite{Radyushkin:2018cvn}是 $\mathcal I(\nu,\mu^2)=\mathfrak M(\nu,z_3^2)+a\int_0^1\dd w\,\mathfrak M(w\nu)\{B(w)[L_z+1]+[4\frac{\ln(1-w)}{1-w}-2(1-w)]_+\}$，在单圈求逆即得 $\sigma=0$ 的 $\mathcal C^{\rm ratio}$；核是实的，$\nu=-\lambda$ 只翻转相位。

(c) 见\cref{ex:lattice-quasi-tree}。

(d) 对夸克靶，由\cref{def:lattice-quasi}，动量空间核是夸克的单圈准 PDF 减去 $f^{(1)}_{q/q}$，也就是 (a)（相应地 (b)）在 $z^2=-\lambda^2/p_z^2$ 的射线上的 Fourier 变换。所需的分布 Fourier 变换是 $\int\frac{\dd\lambda}{2\pi}\ee^{-\ii u\lambda}\ln\lambda^2=-2\gamma_E\delta(u)-[1/|u|]_{\rm reg}$，以及 $\int_0^1\dd\alpha\,g(\alpha)\ee^{\ii\alpha\lambda}$ 型项的变换；结果是\cite[第~III~节]{IzubuchiEtAl2018}所给的\cref{eq:lamet-CMS,eq:lamet-Cratio}。区域 $\xi>1$ 与 $\xi<0$ 上被积函数与\cref{eq:lamet-unphysical}一致；$\pm\frac3{2\xi}$、$\pm\frac3{2(1-\xi)}$ 的移入移出使加号分布内的函数在 $|\xi|\to\infty$ 时可积。我们用 Gauss 型检验函数核对了\cref{eq:lamet-CMS,eq:lamet-Cratio}（两种 $\Gamma$）与 (a)(b) 的 Fourier 变换的配对，二者在数值精度内相同。

(e) $0<\xi<1$：$\frac{1+\xi^2}{1-\xi}=\frac2{1-\xi}-(1+\xi)$，$\ln\xi=\order{1-\xi}$，$\frac{\xi(1+\xi)}{1-\xi}=\frac2{1-\xi}-(2+\xi)$，故方括号中的函数为 $\frac2{1-\xi}(L_\zeta+\ln(1-\xi)-1)+\order{1}$（用 $-\ell_p=L_\zeta$）。$\xi>1$：$\ln\frac\xi{\xi-1}=-\ln(\xi-1)+\order{\xi-1}$，故 $\frac{1+\xi^2}{1-\xi}\ln\frac\xi{\xi-1}=\frac{2\ln(\xi-1)}{\xi-1}+\order{1}$。两侧的奇异项合起来，$\xi<1$ 时 $a(L_\zeta-1)\cdot\frac2{1-\xi}+a\frac{\ln(1-\xi)^2}{1-\xi}$，$\xi>1$ 时 $a(L_\zeta-1)\cdot0+a\frac{\ln(\xi-1)^2}{\xi-1}$，正是所写的统一形式。

(f) $\ell_p=\ln\frac{\mu^2}{4y^2P_z^2}$，$P^z\frac{\dd\ell_p}{\dd P^z}=-2$；\cref{eq:lamet-CMS}只通过 $a\{[\frac{1+\xi^2}{1-\xi}(-\ell_p)]^{[0,1]}_{+(1)}+\frac32\ell_p\delta(1-\xi)\}$ 依赖于 $\ell_p$，故 $P^z\frac{\dd C^{\MSbar}}{\dd P^z}=2a\{[\frac{1+\xi^2}{1-\xi}]_+-\frac32\delta(1-\xi)\}$；代入 $\tilde q=\int\frac{\dd y}{|y|}C\big(\frac xy,\frac\mu{|y|P^z}\big)f(y)$（$f$ 与 $P^z$ 无关）并用 $2a=\as C_F/\pi$ 即得。比值方案中没有 $-\frac32\delta$ 项。
\end{proof}

上述卷积何时收敛？准 PDF 的矩是否存在？（下面 (iii) 中的 $\MSbar$ 条件已见于\cite{IzubuchiEtAl2018}；混合方案与物理夸克分布的部分是本书给出的推导。）

\begin{proposition}[单圈准 PDF 的尾部、矩与卷积定义域 \tagR]\label{prop:lamet-tails}
（$R_{\mathrm{PT}}$：在重整化微扰论中对红外正规化的部分子态严格。）在单圈，记 $C_{\rm fn}$ 为 $\xi\ne1$ 处作为函数的核，则当 $\xi\to\pm\infty$ 时
\begin{equation}\label{eq:lamet-tails}
C^{\MSbar}_{\rm fn}(\xi)=-\frac{3a}{2|\xi|}-\frac{17a}{6}\frac{{\rm sgn}(\xi)}{\xi^2}+\order{|\xi|^{-3}},\qquad C^{\rm ratio}_{\rm fn}(\xi)=-\frac{4a}{3}\frac{{\rm sgn}(\xi)}{\xi^2}+\order{|\xi|^{-3}}.
\end{equation}% VF-LAM-9
因此：
\begin{enumerate}[label=(\roman*),leftmargin=2em,itemsep=1pt,topsep=2pt]
\item $\MSbar$ 准 PDF 在大 $|x|$ 处为 $-\frac{3a}{2|x|}\int f+\order{x^{-2}}$——当 $\int f\neq0$ 时没有零阶矩；这是 $z\to0$ 处 $\ln z^2$ 奇异性的反映。
\item 比值方案中 $\int\dd x\,\tilde q=\int f$ 精确成立（对所有阶，由构造），但 $\tilde q\to-\frac{4a}3\frac{{\rm sgn}(x)}{x^2}\langle y\rangle$，$\langle y\rangle=\int_{-1}^1y\,f(y)\,\dd y=\int_0^1y\,(q+\bar q)\,\dd y$，故一阶矩对数发散，更高阶矩按幂次计数发散。
\item 固定 $x\neq0$，在 $y\to0$ 附近卷积 $\int\frac{\dd y}{|y|}C(\frac xy)f(y)$ 的被积函数为：$\MSbar$，$-\frac{3a}{2|x|}f(y)(1+\order{y})$；混合方案（其参数是 $\lambda_S=z_S|y|P^z$，所以括号在固定 $\lambda_S$ 下的衰减在 $y\to0$ 时没有帮助），$\frac{3a}{2|x|}\big[\frac2\pi{\rm Si}(|x|z_SP^z)-1\big]f(y)(1+\order{y})$；比值方案，$\order{|yf(y)|}$。因此 $\MSbar$ 与混合方案的卷积绝对收敛当且仅当 $\int_{|y|<\epsilon}|f|<\infty$（混合方案除去使方括号为零的孤立的 $|x|z_SP^z$），比值方案只需 $\int|yf|<\infty$。对物理的夸克分布（$y<0$ 处 $f(y)=-\bar q(-y)$），$\order1$ 项合成 $q-\bar q$：只要 $q-\bar q$ 与 $y(q+\bar q)$ 在小 $y$ 处可积，三种方案都给出有限的夸克准 PDF（卷积按对称地挖去 $|y|<\delta$、$\delta\to0^+$ 定义）——包括 Regge 型海夸克 $q,\bar q\sim y^{-\alpha_{\mathbb P}}$、$1\le\alpha_{\mathbb P}<2$ 的情形。胶子核不在本命题范围内。
\end{enumerate}
（本命题只涉及单圈部分子匹配核；$f$ 是任意给定的分布。）
\uses{prop:lamet-oneloop, def:lamet-schemes}
\end{proposition}

\begin{proof}
\emph{第 1 步（化归）。}对固定 $x\neq0$，令 $\xi=x/y$，则 $\frac{\dd y}{|y|}=\frac{\dd\xi}{|\xi|}$，$y\to0$ 对应于 $|\xi|\to\infty$。把核写成 $\delta(1-\xi)c_\delta+\sum_D[g_D]^D_{+(1)}+h$（$h$ 为加号分布之外的函数项）。$\delta$ 项只在 $y=x$ 处起作用；加号分布作用在 $\phi(\xi)=f(x/\xi)/|\xi|$ 上，给出 $\int_Dg_D(\xi)\phi(\xi)\dd\xi-f(x)\int_Dg_D\dd\xi$，后一项是有限常数（\cref{eq:lamet-CMS,eq:lamet-Cratio}中加号分布内的函数在 $|\xi|\to\infty$ 时为 $\order{\xi^{-2}}$，见下）。所以收敛性只取决于 $\int\frac{\dd y}{|y|}C_{\rm fn}(\frac xy)f(y)$ 在 $y\to0$ 附近的行为。

\emph{第 2 步（尾部）。}$\xi>1$：$\ln\frac\xi{\xi-1}=\frac1\xi+\frac1{2\xi^2}+\frac1{3\xi^3}+\cdots$，$\frac{1+\xi^2}{1-\xi}=-(\xi+1+\frac2\xi+\frac2{\xi^2}+\cdots)$，乘积为 $-1-\frac3{2\xi}-\frac{17}{6\xi^2}+\order{\xi^{-3}}$；于是 $\MSbar$ 加号分布内的函数是 $-\frac{17}{6\xi^2}+\order{\xi^{-3}}$，加上外面的 $-\frac3{2\xi}$ 得 $C^{\MSbar}_{\rm fn}=-\frac3{2\xi}-\frac{17}{6\xi^2}+\cdots$。$\xi=-t\to-\infty$：$\frac{1+t^2}{1+t}\ln\frac{1+t}t=1-\frac3{2t}+\frac{17}{6t^2}+\cdots$，加号分布内的函数是 $1-\frac3{2t}+\frac{17}{6t^2}-1+\frac3{2(1+t)}=\frac4{3t^2}+\cdots$，加上外面的 $-\frac3{2(1+t)}$ 得 $-\frac3{2t}+\frac{17}{6t^2}+\cdots$。两侧合起来即\cref{eq:lamet-tails}的第一式。比值核中 $\xi>1$ 的函数为 $-1-\frac3{2\xi}-\frac{17}{6\xi^2}+1+\frac3{2(\xi-1)}=-\frac4{3\xi^2}+\cdots$，$\xi<0$ 为 $+\frac4{3t^2}+\cdots$，即第二式。在坐标空间看：$-\frac{3}{2|\xi|}$ 是 $\frac32\ln\lambda^2\,\delta(1-\alpha)$ 的 Fourier 变换（$\ln\lambda^2\to-1/|u|$），即 $C_0$ 中 $\frac32L_z$ 项的反映；比值核没有 $\delta(1-\alpha)\ln\lambda^2$ 项，其尾部系数 $-\frac43=\int_0^1(\alpha-1)\frac{1+\alpha^2}{1-\alpha}\dd\alpha$ 来自 $-L_zB$ 的一阶矩。

\emph{第 3 步（$y\to0$ 的被积函数）。}由第 2 步，$\frac1{|y|}C_{\rm fn}(\frac xy)$ 中 ${\rm sgn}(\frac xy)\frac{y^2}{|y|}={\rm sgn}(x)\,y$，故 $\MSbar$：$\frac1{|y|}C^{\MSbar}_{\rm fn}=-\frac{3a}{2|x|}-\frac{17a}6\frac{{\rm sgn}(x)}{x^2}y+\order{y^2}$；比值：$-\frac{4a}3\frac{{\rm sgn}(x)}{x^2}y+\order{y^2}$。混合方案：由\cref{eq:lamet-hybrid-conv}与 $C_0=1+a(\frac32L_z+\frac52)$，$|z|>z_S$ 处 $\tilde h_{\rm hyb}-\tilde h_{\rm ratio}=\tilde h_{\rm ratio}[C_0(z)/C_0(z_S)-1]=\frac{3a}2\ln\frac{z^2}{z_S^2}\,\tilde h_{\rm ratio}+\order{a^2}$，即坐标核多出 $\frac{3a}2\delta(1-\alpha)\ln\frac{z^2}{z_S^2}\theta(|z|-z_S)$。沿准 PDF 的射线 $z^2=-\lambda^2/P_z^2$，$\ln\frac{z^2}{z_S^2}=\ln\frac{\lambda^2}{\Lambda_S^2}$，$\Lambda_S\equiv z_SP^z$ \emph{与 $y$ 无关}；因此单圈修正是 $\frac{3a}2\int\dd y\,f(y)\,G(x-y)$，
\[
G(u)=\int\frac{\dd\lambda}{2\pi}\,\ee^{-\ii u\lambda}\ln\frac{\lambda^2}{\Lambda_S^2}\,\theta(|\lambda|-\Lambda_S)=\frac2\pi\int_{\Lambda_S}^\infty\dd\lambda\,\cos(u\lambda)\ln\frac\lambda{\Lambda_S}=-\frac1{|u|}+\frac{2\,{\rm Si}(u\Lambda_S)}{\pi u}\qquad(u\neq0),
\]
最后一步分部积分并用 $\int_0^\infty\frac{\sin(u\lambda)}\lambda\dd\lambda=\frac\pi2{\rm sgn}(u)$；$u=0$ 处的减除由 $\lambda=0$ 处修正为零（$\int G\,\dd u=0$）固定，与 $y\to0$ 无关。即混合方案的修正在强子变量下是平移不变的核 $G(x-y)$（偶函数）。$y\to0$ 时 $G(x-y)\to G(x)=\frac1{|x|}\big[\frac2\pi{\rm Si}(|x|z_SP^z)-1\big]$，一般不为零（它只在 ${\rm Si}(v)=\frac\pi2$ 的孤立点 $v=|x|z_SP^z\approx1.93,\,4.89,\,7.97,\dots$ 为零）。这就是 (iii) 中的三个被积函数。

\emph{第 4 步（收敛性）。}被积函数写成 $k(y)f(y)$，$k$ 在 $y=0$ 的两侧连续。若 $k(0)=c_0\neq0$（$\MSbar$：$c_0=-\frac{3a}{2|x|}$；混合：$c_0=\frac{3a}2G(x)$），则 $kf$ 在零点附近绝对可积当且仅当 $f$ 绝对可积；若 $k(y)=\order{|y|}$（比值），则 $\int|yf|<\infty$ 足够。

\emph{第 5 步（物理夸克分布）。}设 $k(t)=c_0+c_1^+t+\order{t^2}$、$k(-t)=c_0-c_1^-t+\order{t^2}$（$t>0$；两侧斜率可以不同，但 $c_0$ 相同——$\MSbar$ 尾部为偶函数，$G$ 在 $y=0$ 连续）。把 $y=t$ 与 $y=-t$ 的贡献合并，
\[
k(t)q(t)-k(-t)\bar q(t)=c_0\,(q-\bar q)(t)+t\,\big(c_1^+q(t)+c_1^-\bar q(t)\big)+\order{t^2}(|q|+|\bar q|),
\]
它在 $(0,\epsilon)$ 上可积，只要 $q-\bar q$ 与 $t\,q$、$t\,\bar q$ 可积（$q,\bar q\ge0$ 时即 $t(q+\bar q)$ 可积）。于是 $\lim_{\delta\to0^+}\int_{\delta<|y|<\epsilon}$ 存在。对 $q=\bar q\sim y^{-\alpha_{\mathbb P}}$，$y(q+\bar q)\sim y^{1-\alpha_{\mathbb P}}$ 在 $\alpha_{\mathbb P}<2$ 时可积，价分布 $q-\bar q\sim y^{-1/2}$ 可积。这与光锥 ITD $I(\lambda)=\int_0^1\dd t[\cos(t\lambda)(q-\bar q)+\ii\sin(t\lambda)(q+\bar q)]$ 存在的条件相同。胶子与混合核有不同的结构（并作用于 $yg(y)$），这里不作断言。

\emph{第 6 步（(i)(ii)）。}$|x|>1$ 时对所有 $y\in[-1,1]$ 都有 $|\xi|\ge|x|>1$，加号分布的减除项 $\propto f(x)=0$，$\delta(1-\xi)$ 也不起作用，故 $\tilde q^{(1)}(x)=\int\frac{\dd y}{|y|}C_{\rm fn}(\frac xy)f(y)$；第 2 步的展开对 $y$ 一致（余项 $\order{y^2/|x|^3}$），代入得 $\tilde q^{\MSbar}=-\frac{3a}{2|x|}\int f-\frac{17a}6\frac{{\rm sgn}(x)}{x^2}\langle y\rangle+\order{|x|^{-3}}$ 与 $\tilde q^{\rm ratio}=-\frac{4a}3\frac{{\rm sgn}(x)}{x^2}\langle y\rangle+\order{|x|^{-3}}$；$\int_{-1}^1yf\,\dd y=\int_0^1y(q+\bar q)\dd y$ 来自 $f(-t)=-\bar q(t)$。于是 $\MSbar$ 的零阶矩、比值方案的一阶矩都以对数方式发散。比值方案中 $\int\dd x\,\tilde q=\int\dd\xi\,C^{\rm ratio}(\xi)\int\dd y\,f(y)$（Fubini，此时一切绝对收敛），而单圈的修正全是加号分布，$\int C^{\rm ratio}=1$；对所有阶，$\mathfrak M(0,z^2)=1$ 恒成立，故 $\int\tilde q=1$（按价夸克数归一化）。
\end{proof}

用部分子变量 $\xi=x/y$，$G(x-y)=\frac1{|y|}\mathcal B(\xi)$，$\mathcal B(\xi)=-\frac1{|1-\xi|}+\frac{2\,{\rm Si}((1-\xi)\lambda_S)}{\pi(1-\xi)}$，$\lambda_S=z_S|y|P^z$，正是下一节混合核中的括号；$\lambda_S\to0$ 时括号的 $1/u^2$ 衰减消失，所以混合方案与 $\MSbar$ 需要同样的小 $y$ 条件。

\begin{proposition}[非单态准 PDF 匹配核的两圈结果 \tagR]\label{prop:lamet-twoloop}
（$R_{\mathrm{PT}}$：在重整化微扰论中对红外正规化的部分子态严格。）在两圈，非单态坐标空间系数函数以及相应的准 PDF 与赝 PDF 匹配核已由两个小组独立算出\cite{Li:2020xml,Chen:2020ody}（味非对角部分\cite{Chen:2020arf}，主积分\cite{Chen:2020iqi}；据其摘要）；按文献\cite{Chen:2020ody}摘要的说法，这一计算“在这一阶显式地建立了非平凡的因子化”。两圈核满足重整化群方程，并与第~23~章的阈值因子化在两圈相容。
（本命题只涉及部分子态。）
\uses{prop:lamet-oneloop}
\end{proposition}

\begin{proof}
这是文献中的显式计算\cite{Li:2020xml,Chen:2020ody}，两组结果一致；公式不在此重复。重整化群相容性：两圈核的 $\ln\mu^2$ 依赖由两圈 DGLAP 核与两圈算符反常维度决定（后者见\cite{Braun:2020ymy}）。与阈值因子化的比较见第~23~章；那里有两个常数是由与两圈结果的比较定出的，所以这一比较是相容性检验，不是独立验证。
\end{proof}

更高阶：非单态 Wilson 线夸克关联函数的完整 N$^3$LO 匹配也已算出\cite{Cheng:2024wyu}（文献\cite{Ji:2024hit}把它用作 Mellin 空间中的三圈结果）；三圈反常维度已知\cite{Braun:2020ymy}；带重求和的 NNLO 匹配已用于格点分析\cite{Gao:2021dbh,Su:2022fiu}。

\begin{proposition}[极化、横度与扭度 3 准分布的匹配 \tagP]\label{prop:lamet-polarized}
螺旋度（$\gamma^z\gamma_5$）与横度（$\gamma^0\gamma^\perp\gamma_5$）准 PDF 在单圈与相应的光锥分布匹配（截断方案\cite{XiongEtAl2014}；RI/MOM 方案\cite{Lin:2018pvv,Liu:2018hxv}，后者为预印本；据其摘要）。扭度 3：$g_T$\cite{Bhattacharya:2020xlt}；$e$ 与 $h_L$，其中有正比于 $\delta(x)$ 的本质零模贡献\cite{Bhattacharya:2020jfj}；轴矢\cite{Braun:2021aon}与手征奇\cite{Braun:2021gvv}准、赝分布中扭度 2 与扭度 3 的分离借助 Jaffe--Ji 型关系（均据其摘要）。
（本命题只涉及部分子态。）
\uses{def:lattice-quasi, def:pdf-polarized, prop:lattice-renormalizability}
\end{proposition}

\paragraph{地位与缺口}
\tagP：只有单圈匹配；包括零模在内的扭度 3 全阶匹配尚未建立；以上引述均据各文摘要（\cite{Liu:2018hxv}为预印本）；$\gamma_5$ 的方案按\cref{def:prelim-gamma5}处理。一般的陈述——极化准分布满足与\cref{thm:lattice-lamet}同样的全阶匹配——是自然的推广，但没有在文献中写出。
